\subsection{Transposition of a continuous linear mapping}

In this No., we suppose that \((\mathrm{F},\mathrm{G})\) and \((\mathrm{F}_1,\mathrm{G}_1)\) are two vector spaces in duality.

PROPOSITION 5. --- Let u be a linear mapping of F in \(F_{1}\) . The following properties are equivalent :

\begin{enumerate}
\def\labelenumi{\alph{enumi})}
\item
  \(u\) is continuous for the weak topologies \(\sigma (\mathrm{F},\mathrm{G})\) and \(\sigma (\mathrm{F}_1,\mathrm{G}_1)\) ;
\item
  there exists a mapping \(v: \mathbf{G}_1 \to \mathbf{G}\) such that
\end{enumerate}

\[
\langle u (y), z _ {1} \rangle = \langle y, v (z _ {1}) \rangle\tag{1}
\]

for all \(y \in \mathbf{F}\) and \(z \in \mathbf{G}_1\) .

If these properties hold and if the duality between \(\mathbf{F}\) and \(\mathbf{G}\) is separating in \(\mathbf{G}\) , then the mapping \(v\) satisfying (1) is unique, and \(v\) is linear.

If u is continuous for the weak topologies, then, for all \(z_{1} \in G_{1}\) , the linear form \(y \mapsto \langle u(y), z_{1} \rangle\) on F is continuous for \(\sigma(F, G)\) , thus (II, p.~43, prop. 3) can be written as \(y \mapsto \langle y, v(z_{1}) \rangle\) with \(v(z_{1}) \in G\) , which shows that a) implies b). Conversely, if b) is true, for all \(z_{1} \in G_{1}\) , the linear form

\[
y \mapsto \langle y, v (z _ {1}) \rangle = \langle u (y), z _ {1} \rangle
\]

is continuous for \(\sigma(\mathrm{F},\mathrm{G})\) : it follows from the definition of weak topologies that \(u\) is continuous for \(\sigma(\mathrm{F},\mathrm{G})\) and \(\sigma(\mathrm{F}_1,\mathrm{G}_1)\) (I, p.~10, cor. 1). The uniqueness of \(v\) follows from \((\mathrm{D}_{\mathrm{II}})\) and this uniqueness implies that \(v\) is linear.

Remark. --- Suppose that the duality between F and G is separating in G and that the duality between \(F_{1}\) and \(G_{1}\) is separating in \(G_{1}\) . If we identify G and \(G_{1}\) with subspaces of \(F^{*}\) and \(F_{1}^{*}\) respectively, the conditions a) and b) are equivalent to \(^{t}u(\mathbf{G}_{1}) \subset \mathbf{G}\) ; v is the restriction of the transpose \(^{t}u\) of u (A, II, §2.5) to \(G_{1}\) .

We say, simply (when there is no chance of confusion) that v is the transpose of u (relative to the duality on the one hand between F and G and on the other hand between \(F_{1}\) and \(G_{1}\) ) and we again use \(^{t}u\) to denote it.

COROLLARY. --- Suppose that the duality between F and G is separating in G. If u is a linear mapping of F in \(F_{1}\) , that is continuous for \(\sigma(F, G)\) and \(\sigma(F_{1}, G_{1})\) , then its transpose is a linear mapping of \(\mathbf{G}_1\) in \(\mathbf{G}\) , that is continuous for \(\sigma (\mathbf{G}_1,\mathbf{F}_1)\) and \(\sigma (\mathbf{G},\mathbf{F})\) . Further if the duality between \(\mathbf{F}_1\) and \(\mathbf{G}_1\) is separating in \(\mathbf{F}_1\) then \(^{t}(^{t} u) = u\) .

It is sufficient to exchange F and \(F_{1}\) with G and \(G_{1}\) in prop. 5.

PROPOSITION 6. --- Suppose that the duality between F and G (resp. \(F_{1}\) and \(G_{1}\) ) is separating in G (resp. \(F_{1}\) ). Let u be a linear mapping of F in \(F_{1}\) that is continuous for \(\sigma(F, G)\) and \(\sigma(F_{1}, G_{1})\) . Let A be a set in F and \(A_{1}\) a set in \(F_{1}\) ; then :

\begin{enumerate}
\def\labelenumi{(\roman{enumi})}
\item
  We have \((u(\mathbf{A}))^{\circ} = ({}^{t} u)^{-1} (\mathbf{A}^{\circ})\) .
\item
  We have \(\overline{{}^{t}u(A_1^\circ)}\subset (u^{-1}(A_1))^{\circ}\) ; further, if \(A_1\) is closed, (for \(\sigma(F_1,G_1)\) ) convex, and contains the origin, then we have \(\overline{{}^{t}u(A_1^\circ)} = (u^{-1}(A_1))^{\circ}\) .
\end{enumerate}

Let \(z_{1} \in G_{1}\) , the relation \(z_{1} \in (u(A))^{\circ}\) is equivalent to \(\langle u(y), z_{1} \rangle \geqslant -1\) for all \(y \in A\) , and the relation \({}^{t}u(z_{1}) \in A^{\circ}\) is equivalent to \(\langle y, {}^{t}u(z_{1}) \rangle \geqslant -1\) for all \(y \in A\) and our assertion (i) follows using (1). Next interchanging \(u\) and \({}^{t}u\) and applying (i) to the set \(A_{1}^{\circ}\) of \(G_{1}\) we get

\[
\left(^ {t} u \left(\mathrm{A} _ {1} ^ {\circ}\right)\right) ^ {\circ} = u ^ {- 1} \left(\mathrm{A} _ {1} ^ {\circ \circ}\right) \supset u ^ {- 1} \left(\mathrm{A} _ {1}\right)\tag{2}
\]

from which, on taking polars

\[
\left(^ {t} u \left(\mathrm{A} _ {1} ^ {\circ}\right)\right) ^ {\circ \circ} \subset \left(u ^ {- 1} \left(\mathrm{A} _ {1}\right)\right) ^ {\circ}.
\]

We have \(\overline{(^t u(A_1^\circ))} \subset (^{t}u(A_1^\circ))^{\circ \circ}\) by the bipolar theorem (II, p.~44, th. 1); the final statement follows from (2) and the bipolar theorem since then \(A_1^{\circ \circ} = A_1\) and \(^t u(A_1^\circ)\) is convex and contains the origin.

COROLLARY 1. --- With the notations of prop. 6, the relation \(u(\mathbf{A}) \subset \mathbf{A}_1\) implies \(^t u(\mathbf{A}_1^\circ) \subset \mathbf{A}^\circ\) ; if further \(\mathbf{A}_1\) is convex, closed (for \(\sigma(\mathbf{F}_1, \mathbf{G}_1)\) ) and contains the origin, then these two relations are equivalent.

In fact, the relation \(u(\mathbf{A}) \subset \mathbf{A}_1\) is equivalent to \(\mathbf{A} \subset u^{-1}(\mathbf{A}_1)\) , therefore implies

\[
{ } ^ { t } u ( \mathrm{A} _ { 1 } ^ { \circ } ) \subset \overline { { { ^ { t } u ( \mathrm{A} _ { 1 } ^ { \circ } ) } } } \subset ( u ^ { - 1 } ( \mathrm{A} _ { 1 } ) ) ^ { \circ } \subset \mathrm{A} ^ { \circ }
\]

and conversely the relation \({}^t u(\mathbf{A}_1^\circ) \subset \mathbf{A}^\circ\) implies

\[
\mathrm{A} ^ {\circ \circ} \subset \left(^ {t} u (\mathrm{A} _ {1} ^ {\circ})\right) ^ {\circ} = u ^ {- 1} (\mathrm{A} _ {1} ^ {\circ \circ})
\]

from (2). When \(\mathrm{A}_1 = \mathrm{A}_1^{\circ \circ}\) , we deduce that \(\mathrm{A} \subset u^{-1}(\mathrm{A}_1)\) .

COROLLARY 2. --- Let u be a linear mapping of F in \(F_{1}\) that is continuous for \(\sigma(F, G)\) and \(\sigma(F_{1}, G_{1})\) . We have then

\[
\operatorname{Ker} (^ {t} u) = (\operatorname{Im} (u)) ^ {\circ},\tag{3}
\]

\[
\overline {{\operatorname{Im} (^ {t} u)}} = (\operatorname{Ker} (u)) ^ {\circ}.\tag{4}
\]

Suppose that the dualities between F and G and between \(F_{1}\) and \(G_{1}\) are separating; then \(u(F)\) is dense in \(F_{1}\) (for \(\sigma(F_{1}, G_{1})\) ), if and only if \(^{t}u\) is injective.

Apply prop. 6 with A = F and \(A_{1} = \{0\}\) , using the fact that the weak topologies \(\sigma(G, F)\) and \(\sigma(F_{1}, G_{1})\) are Hausdorff. The last assertion results from (4), interchanging u and \(^{t}u\) .
% semantic-fragments: legacy-input-seam
\relax{}
